\section{Related work}

\subsection{The benchmark and its two sanctioned protocols}

JIGSAWS \cite{gao2014jigsaws} is the standard public benchmark for robot-assisted surgical skill assessment, and the accompanying benchmark suite \cite{ahmidi2017dataset} defines the two cross-validation schemes used throughout this literature. Leave-one-supertrial-out holds out the \textit{i}-th repetition across all surgeons; leave-one-user-out holds out every trial from one surgeon. Both ship with the dataset and both are sanctioned, so the protocol we examine is the benchmark's own rather than an informal practice.

The field is aware that individual style exists. A review of gesture recognition in robotic surgery \cite{vanamsterdam2021gesture} notes that leave-one-user-out is the harder setting ``due to subject-specific style variability'', and surveys methods that normalise such variability away. That framing treats identity as nuisance variance to be suppressed in pursuit of a better model. Our concern is the reverse: the same variability, under a protocol that places a surgeon on both sides of the split, does not degrade a score but inflates it.

\subsection{Label validity and protocol validity are separate critiques}

The closest existing critique of skill assessment on this benchmark concerns the target. Hendricks et al. \cite{hendricks2024exploring} show that the self-declared novice/intermediate/expert label correlates poorly with the expert GRS rating, and \cite{hoffmann2026looking} reports the same mismatch. Our own data reproduce it --- one surgeon labelled expert averages GRS 13.5 while one labelled novice averages 19.1 (Section~3.1) --- and we report that as a replication of their finding rather than as one of ours.

The distinction from our contribution is worth stating precisely, because the two are easily conflated. \textbf{Hendricks et al. ask whether the target measures skill. We ask whether the evaluation protocol permits a model to succeed without measuring skill at all, whatever the target's quality.} These critiques are independent, and they compound. A perfectly valid target would not repair what we describe: surgeon identifiability is a property of the feature representation and the fold structure, not of what the label means, so a model could shortcut a flawless label just as readily. Conversely their critique survives a flawless protocol, since a well-designed split cannot rescue a target that measures the wrong thing. A reader persuaded by Hendricks et al. should not conclude that fixing the label would fix the benchmark.

\subsection{The performance gap is widely reported, and framed three other ways}

That accuracy falls under leave-one-user-out is not news in this literature. What differs is the explanation offered.

An interpretable whole-stream video framework \cite{chen2026interpretable} reports near-perfect supertrial-out accuracy across the three tasks, a significant drop under user-out evaluation, and near-chance transfer to an external dataset, concluding that generalisation to unseen surgeons is limited and calling for ``surgeon-invariant feature learning''; the goal is named, the mechanism is not measured. \cite{hoffmann2026looking} evaluates GRS prediction under a strict disjoint-participant protocol and reports a concordance correlation of -0.03 (+/-0.08) for a supervised model trained and evaluated entirely on JIGSAWS, below chance, against 0.22 for its best baselines; having ruled out model-side causes through an extensive architecture and hyperparameter search, it attributes the failure to the annotations, which it reports ``lack a consistent ordering with respect to the visual features that encode actual skill differences''. A recent explainable-machine-learning study using rich kinematic features \cite{xuan2026explainable} criticises prior JIGSAWS work for lacking subject-independent validation --- the same complaint, again without a direct test. Lefor et al. \cite{lefor2020motion} earlier flagged the dataset's subject imbalance and an inverse correlation among subjects as a threat to generalisation, without quantifying either.

Closest to our own framing is the markerless hand-tracking study of open-surgery suturing by Constans et al. \cite{constans2026markerless}, which finds accuracy falling to 52.3\% under leave-one-subject-out and states that this gap ``quantifies the surgeon-specific component captured by video-level splits''. We should be plain about how little separates that sentence from our thesis: they articulate the same interpretation of the same kind of gap, and had they run a subject-identification probe, the attribution would already exist in this literature. Our contribution relative to that study is correspondingly narrower than relative to the others above --- a different modality, robot kinematics rather than open-surgery hand tracking, and a direct measurement of the component they name but do not test. Taken together these studies establish that the field suspects a subject-specific confound; none tests it directly.

The null in ``Looking Beyond the Scale'' overlaps our Section~4.3, and that overlap strengthens rather than duplicates the contribution. Their study and ours reach the same conclusion about cross-participant GRS prediction through almost entirely different routes: different feature representations, different model families, a different agreement metric, and a different reference against which the result is judged. A single null is always vulnerable to the objection that it reflects an idiosyncrasy of one pipeline --- an underpowered search, an unlucky architecture, a poorly chosen metric. Two independent pipelines arriving at the same place is much harder to explain that way, and their architecture search excludes model-side causes that our single linear model could not have excluded alone. Independent corroboration by a different method is better evidence than either result on its own. What their study does not supply, and what we add, is the attribution: they establish that the signal is absent, while we measure where the apparent signal under the permissive protocol comes from.

One counter-argument must be met directly, because it applies to every study above and to our Section~4.1. Responding to Saeb et al. \cite{saeb2017need}, Little et al. \cite{little2017using} argued that a gap between record-wise and subject-wise cross-validation can arise from \textit{underfitting} on heterogeneous data rather than from identity confounding, so that a gap alone establishes nothing about mechanism. The objection is sound, and none of the surgical studies above can answer it, because each reports only a gap. Our Section~4.2 is designed to answer it: the same 16 features, the same model and the same folds identify which of eight surgeons performed a held-out repetition with accuracy between 0.83 and 0.93, against chance rates near 0.13. A performance gap is ambiguous between confounding and underfitting; recoverable identity at roughly seven times chance is not.

Surgeon identifiability itself is not a new observation. Working from intraoperative tool-tissue force data rather than benchmark kinematics, Baghdadi et al. \cite{baghdadi2023force} identify individual surgeons with micro-average AUC 0.97 and 83.8\% accuracy across 89 neurosurgical cases, concluding that each surgeon has a unique force profile recoverable by deep networks. They frame that finding as a capability: a route to bespoke per-surgeon feedback and to correlating outcomes with individual performance. It is not connected there to the validity of any skill-assessment protocol, and no performance gap is reported.

Our claim is correspondingly narrow. We are not the first to identify surgeons from surgical signals; what we add is the use of that identifiability as a validity diagnostic, run on the same features, model and folds as the skill task, and tied to the protocol gap it explains. To our knowledge no study both attributes that gap to surgeon-identity leakage and demonstrates the leakage by classifying the surgeon from the features used for skill. This is the outcome of a recorded search of PubMed and general web sources (\texttt{paper/systematic\_search.md}); Google Scholar and IEEE Xplore could not be queried programmatically, so the claim is bounded rather than exhaustive.

\subsection{Reported performance, and the missing reference}

In the studies we examined \cite{zia2018automated,wang2021towards,vbanet2022video}, results are reported as accuracies or as correlations between predicted and expert scores, without the value the same metric would take for a predictor using no motion information at all. We have not surveyed this literature systematically enough to call that a field-wide norm; it is the pattern in the work we read. We depart from it deliberately, for the reason given in Section~3.5 and Section~3.7: under leave-one-user-out, a predictor that returns its training fold's mean GRS scores approximately -0.8 on pooled rank correlation purely by construction. A positive correlation is therefore not self-interpreting, and a model can post one while committing larger absolute errors than that uninformative predictor.

\subsection{Calibration and selective prediction: a borrowed lineage}

Calibration has begun to receive attention in surgery: \cite{surgiq2026} benchmarks large language models on surgical knowledge questions and includes a calibration analysis, reporting that even strong models make confident mistakes on plausible distractors. That is a different setting from ours. We located no calibration or selective-prediction study of a surgical \textit{skill assessment} model built from motion, kinematic or video data, and surgical work on confidence-gated deferral is correspondingly thin. The methodological lineage we draw on is therefore from diagnostic imaging and intensive care, where deferral is established and demonstrably useful: CoDoC \cite{dvijotham2023codoc} reduces false positives by 25\% at matched sensitivity while cutting clinician workload, and ``Trust or Escalate'' \cite{zheng2024trust} provides calibrated selective evaluation with provable agreement guarantees. Overconfidence in model-as-judge settings is itself an active topic \cite{tian2025overconfidence}. We take from this literature one caution that bears directly on our own null result: deferral cannot manufacture reliability that the underlying model lacks \cite{calibratedtriage2026}.

\subsection{Zero-shot and text-based scoring}

Foundation models have reached this benchmark, but not in the role we examine. The nearest work, ZEAL \cite{zeal}, uses zero-shot video segmentation to produce features for a \textit{trained} LSTM, so the learning still happens downstream; that is a different mechanism from a decision model scoring a rubric with no training examples at all. The model we evaluate, Jev, is marketed on calibration --- its training method is branded Reinforcement Learning for Calibrated Decisions --- though a trade critique notes the term is coined and that its ``zero hallucinations'' claim concerns schema conformance rather than correctness \cite{eliot2026forbes}. Rather than rely on vendor claims, we note three independent evaluations published in September 2026 \cite{priorbench2026,scienthoon2026,webofmike2026} converging on a consistent pattern: calibration holds on in-distribution data and degrades markedly out of distribution, and one reports that the model's separate confidence field is less informative than its maximum class probability. We give no figures from these evaluations, for the reason that follows. \textbf{All three are non-peer-reviewed single-author repositories or blog posts, and we cite them as preliminary external evidence rather than settled results}; they also identify the model by version strings that may not correspond exactly to the build we pinned, so we treat them as evidence about the model family rather than as replications of our setup. The pattern they describe is the one our Section~4.5 reports on unseen surgeons.
